% !TeX root = ../../Book3.tex
\section{素数定理}
\label{b3:sec:2705}

Tauber 步骤给出的是素数幂的对数权和。现在把这个结果代入本章开头已经证明的等价关系，就能得到通常叙述的素数定理。本节还说明它允许怎样改变计数方式，以及哪些看似相近的定量结论仍未由此证明。

\subsection{素数定理及其加权推论}
\label{b3:subsec:270501}

\begin{theorem}[素数定理]\label{b3:thm:27-pnt}
对实数 \(x\ge2\)，令 \(\pi(x)=\#\{\,p\le x\mid p\text{ 为素数}\,\}\)、\(\vartheta(x)=\sum_{p\le x}\log p\)、\(\psi(x)=\sum_{p^m\le x,\ m\ge1}\log p\)，其中 \(p\) 遍历素数、\(m\) 为正整数。则当 \(x\to\infty\) 时，
\[
 \psi(x)\sim x,
 \quad \vartheta(x)\sim x,
 \quad \pi(x)\sim\frac{x}{\log x}.
\]
\end{theorem}
\begin{proof}
在\subsecref{b3:subsec:270403}中，已经核对了 \(A(x)=\psi(x)\) 右连续、非负、非递减，其 Stieltjes 变换在 \(\operatorname{Re}s>1\) 等于 \(-\zeta'(s)/\zeta(s)\)，且减去 \(1/(s-1)\) 后全纯越过整条边界直线。\cref{b3:thm:27-wiener-ikehara}取 \(c=1\)，给出 \(\psi(x)/x\to1\)。再应用\cref{b3:prop:27-equivalences}，就得到其余两个渐近式；其中去掉高次素数幂只损失 \(O(\sqrt x\log x)=o(x)\)，去掉对数权的积分余项也已在该命题中估计。
\end{proof}

这个结果称为\term[sushu-dingli]{素数定理}[prime number theorem]。它断言相对误差趋零，却没有给出趋零的数值速率。即使知道 \(\pi(x)\) 的整体主项，也不能把两个很接近的端点相减后，仍无条件保留同样的相对精度。

\begin{corollary}[固定比例区间内的素数]
\label{b3:cor:27-fixed-ratio-interval}
设 \(\lambda>1\) 为固定实数，\(\pi(x)\) 为不超过实数 \(x\ge1\) 的素数个数。则当 \(x\to\infty\) 时，
\[
 \#\{\,p\mid x<p\le\lambda x,\ p\text{ 为素数}\,\}
 =\pi(\lambda x)-\pi(x)
 \sim(\lambda-1)\frac{x}{\log x}.
\]
特别地，对这个固定的 \(\lambda\)，存在 \(X_\lambda\)，使每个实数 \(x\ge X_\lambda\) 的区间 \((x,\lambda x]\) 都含有素数。
\end{corollary}
\begin{proof}
素数定理给出
\[
 \begin{aligned}
 \pi(\lambda x)-\pi(x)
 &=\frac{\lambda x}{\log(\lambda x)}-\frac{x}{\log x}
      +o\!\left(\frac{x}{\log x}\right)\\
 &=\frac{x}{\log x}
    \left(\frac{\lambda\log x}{\log x+\log\lambda}-1+o(1)\right).
 \end{aligned}
\]
因为 \(\lambda\) 固定，括号趋于正数 \(\lambda-1\)，证明了渐近式。右侧趋于无穷，所以区间内的素数个数最终为正。
\end{proof}

若让 \(\lambda=\lambda(x)\downarrow1\)，预期主项会缩小到 \((\lambda(x)-1)x/\log x\)，而我们目前只有两个端点误差各为 \(o(x/\log x)\)。这不足以保证误差比新的主项小。因此这里的“固定比例”是推论实际使用的条件；任意短区间的计数还需要额外估计。

\ProseParagraphBreak

去掉或加入一个平滑权重仍由部分求和完成。下面把主项中的常数也算出，以免把“乘上 \(p^\alpha\)”误解为简单地把 \(\pi(x)\) 乘上 \(x^\alpha\)。

\begin{proposition}[素数的幂权和]
\label{b3:prop:27-prime-power-weights}
设 \(\alpha>-1\) 为固定实数。对于实数 \(x\ge2\)，对不超过 \(x\) 的素数求和，则
\[
 \sum_{p\le x}p^\alpha
 \sim\frac{x^{\alpha+1}}{(\alpha+1)\log x}
 \quad(x\to\infty).
\]
渐近关系中的参数 \(\alpha\) 保持固定，不主张对趋近于 \(-1\) 的参数一致成立。
\end{proposition}
\begin{proof}
先计算所需的积分尺度。置 \(\beta=\alpha+1>0\)，并令
\[
 I_\alpha(x)=\int_2^x\frac{t^\alpha}{\log t}\,\mathrm dt,
 \quad D_\alpha(x)=\frac{x^\beta}{\beta\log x}.
\]
这里 \(D_\alpha(x)\to\infty\)，且
\[
 D_\alpha'(x)=\frac{x^\alpha}{\log x}
       \left(1-\frac1{\beta\log x}\right).
\]
给定 \(0<\varepsilon<1\)，当 \(x\ge X\) 且 \(X\) 足够大时，
\((1-\varepsilon)I_\alpha'(x)\le D_\alpha'(x)\le I_\alpha'(x)\)。从 \(X\) 到 \(x\) 积分，固定下端的常数除以 \(D_\alpha(x)\) 后消失，于是
\[
 1\le\liminf_{x\to\infty}\frac{I_\alpha(x)}{D_\alpha(x)}
 \le\limsup_{x\to\infty}\frac{I_\alpha(x)}{D_\alpha(x)}
 \le\frac1{1-\varepsilon}.
\]
令 \(\varepsilon\downarrow0\)，得到 \(I_\alpha(x)\sim D_\alpha(x)\)。这个论证同时适用于正的和负的 \(\alpha\)，只使用了 \(\beta>0\)。

记 \(r(t)=\pi(t)-t/\log t=o(t/\log t)\)。由\cref{b3:prop:25-integral-summation}，
\[
 \sum_{p\le x}p^\alpha
 =x^\alpha\pi(x)-\alpha\int_2^x t^{\alpha-1}\pi(t)\,\mathrm dt.
\]
边界项中的误差 \(x^\alpha r(x)\) 为 \(o(x^\beta/\log x)\)。对积分误差，选 \(X\) 使 \(|r(t)|\le\varepsilon t/\log t\) 对所有 \(t\ge X\) 成立，则
\[
 \left|\int_2^x t^{\alpha-1}r(t)\,\mathrm dt\right|
 \le C_{\alpha,X}+\varepsilon I_\alpha(x).
\]
除以 \(D_\alpha(x)\)，先令 \(x\to\infty\)，再令 \(\varepsilon\downarrow0\)，可知该误差为 \(o(D_\alpha(x))\)。这里初始区间只给固定常数；\(D_\alpha(x)\to\infty\) 正是它能够被忽略的原因。因此
\[
 \begin{aligned}
 \sum_{p\le x}p^\alpha
 &=\frac{x^\beta}{\log x}-\alpha I_\alpha(x)
           +o\!\left(\frac{x^\beta}{\log x}\right)\\
 &=\left(1-\frac\alpha{\alpha+1}\right)
      \frac{x^\beta}{\log x}
           +o\!\left(\frac{x^\beta}{\log x}\right),
 \end{aligned}
\]
即为所求。
\end{proof}

例如，取 \(\alpha=1\) 得 \(\sum_{p\le x}p\sim x^2/(2\log x)\)；取 \(\alpha=-1/2\) 得 \(\sum_{p\le x}p^{-1/2}\sim2\sqrt x/\log x\)。第二个例子也说明，权重虽在衰减，累积和仍可增长。临界值 \(\alpha=-1\) 则必须单独处理，不能把它代入分母 \(\alpha+1\)。

\begin{proposition}[素数倒数和的主项]
\label{b3:prop:27-prime-reciprocal-asymptotic}
对于实数 \(x\ge2\)，对不超过 \(x\) 的素数求和，则当 \(x\to\infty\) 时，
\[
 \sum_{p\le x}\frac1p=\log\log x+o(\log\log x).
\]
\end{proposition}
\begin{proof}
部分求和给出
\[
 \sum_{p\le x}\frac1p
 =\frac{\pi(x)}x+\int_2^x\frac{\pi(t)}{t^2}\,\mathrm dt.
\]
边界项为 \(O(1/\log x)\)。写 \(\pi(t)=t/\log t+r(t)\)，其中 \(r(t)=o(t/\log t)\)，则积分主项为
\[
 \int_2^x\frac{\mathrm dt}{t\log t}
 =\log\log x-\log\log2.
\]
给定 \(\varepsilon>0\)，取固定 \(X>2\)，使 \(|r(t)|\le\varepsilon t/\log t\) 对所有 \(t\ge X\) 成立。于是
\[
 \left|\int_2^x\frac{r(t)}{t^2}\,\mathrm dt\right|
 \le C_X+\varepsilon(\log\log x-\log\log X).
\]
除以 \(\log\log x\)，先令 \(x\to\infty\)，再令 \(\varepsilon\downarrow0\)，得到所需余项。
\end{proof}

这个结果给出了素数倒数和发散的主项。误差 \(o(\log\log x)\) 本身仍可能无界；要判断 \(\sum_{p\le x}1/p-\log\log x\) 是否收敛到常数，需要进一步控制误差积分。

\ProseParagraphBreak

素数定理还有两个常用的表述：给定素数的序号，估计这个素数有多大；或者用对数积分作为比较函数。两者都是主项层面的改写，下面分别说明如何恢复其中的对数和误差尺度。

\begin{corollary}[第 \texorpdfstring{\(n\)}{n} 个素数的渐近]
\label{b3:cor:27-nth-prime}
设 \(p_n\) 为按递增顺序排列的第 \(n\) 个素数，\(n\ge1\)。则
\[
 p_n\sim n\log n\quad(n\to\infty).
\]
\end{corollary}
\begin{proof}
素数定理保证 \(\pi(x)\to\infty\)，所以每个 \(p_n\) 均存在，且 \(p_n\to\infty\)。在 \(x=p_n\) 处有精确等式 \(\pi(p_n)=n\)，因而
\[
 n=\frac{p_n}{\log p_n}(1+o(1)).
\]
这里不能先把 \(\log p_n\) 擅自换成 \(\log n\)，因为这正是尚须证明的一步。对上式取实对数，由 \(\log(1+o(1))=o(1)\)，有
\[
 \log n=\log p_n-\log\log p_n+o(1).
\]
除以趋于无穷的 \(\log p_n\)，并使用 \(\log\log p_n/\log p_n\to0\)，得到 \(\log n/\log p_n\to1\)。最后把两个已知比值相乘：
\[
 \frac{p_n}{n\log n}
 =\frac{p_n}{n\log p_n}\cdot\frac{\log p_n}{\log n}
 \longrightarrow1.
\]
\end{proof}

这个推导展示了反求渐近量时的一般步骤：先从计数关系取得含未知量的等式，再证明缓慢变化的对数可以替换，最后代回。它没有把“函数近似于某函数”直接当作“反函数近似于随意代换后的表达式”。

\begin{proposition}[对数积分与主项尺度]
\label{b3:prop:27-logarithmic-integral}
对实数 \(x\ge2\)，定义
\[
 \operatorname{Li}(x)\coloneqq\int_2^x\frac{\mathrm dt}{\log t}.
\]
则当 \(x\to\infty\) 时，
\[
 \operatorname{Li}(x)=\frac{x}{\log x}
          +O\!\left(\frac{x}{\log^2x}\right),
 \quad
 \pi(x)\sim\operatorname{Li}(x),
\]
其中 \(\pi(x)\) 是不超过 \(x\) 的素数个数。
\end{proposition}
\begin{proof}
由 \((t/\log t)'=1/\log t-1/\log^2t\)，分部积分给
\[
 \operatorname{Li}(x)
 =\frac{x}{\log x}-\frac2{\log2}
             +\int_2^x\frac{\mathrm dt}{\log^2t}.
\]
当 \(x\ge4\) 时，把余项分在 \(\sqrt x\)，有
\[
 0\le\int_2^x\frac{\mathrm dt}{\log^2t}
 \le\frac{\sqrt x}{\log^22}+\frac{4x}{\log^2x}
 =O\!\left(\frac{x}{\log^2x}\right).
\]
最后一步使用 \(\log^2x/\sqrt x\to0\)。于是 \(\operatorname{Li}(x)\sim x/\log x\)；再与素数定理比较，得到 \(\pi(x)/\operatorname{Li}(x)\to1\)。这个比值只在足够大的 \(x\) 上使用，因而 \(\operatorname{Li}(2)=0\) 不造成问题。
\end{proof}

这里的\term[duishu-jifen]{对数积分}[logarithmic integral]从二开始积分，不涉及一处的奇点；采用其他常用归一化时可能相差一个常数，本章固定使用上面的定义。由当前证明，我们只可断言
\[
 \pi(x)-\operatorname{Li}(x)=o\!\left(\frac{x}{\log x}\right).
\]
\(\operatorname{Li}(x)\) 与 \(x/\log x\) 之间的差有更小的上界，并不意味着 \(\pi(x)-\operatorname{Li}(x)\) 自动也有这个上界。要证明对数积分提供更精细的误差近似，必须另行估计素数计数的余项。

\subsection{误差项与零点的关系}
\label{b3:subsec:270502}

素数定理给出 \(\psi(x)=x+o(x)\)，它刻画了误差相对于主项的大小。若误差能够改进为一个严格小于一的固定幂次，变换就会向左多收敛一段；由此可以严格推出一个无零点半平面。误差积分的绝对收敛将使这一联系具体化。

\begin{proposition}[幂次误差蕴含无零点半平面]
\label{b3:prop:27-power-error-zero-free}
设 \(\zeta\) 为 Riemann zeta 函数的亚纯延拓，且对于实数 \(x\ge1\) 定义 \(\psi(x)=\sum_{p^r\le x,\ r\ge1}\log p\)，其中 \(p\) 遍历素数、\(r\) 为正整数。若存在固定实数 \(0\le\alpha<1\)、\(B\ge0\)，使
\[
 \psi(x)=x+O\bigl(x^\alpha(\log x)^B\bigr)
 \quad(x\to\infty),
\]
则 \(\zeta(s)\) 在严格半平面 \(\operatorname{Re}s>\alpha\) 内没有零点；\(s=1\) 仍为单极点。
\end{proposition}
\begin{proof}
令 \(E(x)=\psi(x)-x\)。由假设并增大常数，可对所有 \(x\ge2\) 使用 \(|E(x)|\le Cx^\alpha(\log x)^B\)。在 \(\operatorname{Re}s>\alpha\) 定义
\[
 H(s)=s\int_1^\infty E(x)x^{-s-1}\,\mathrm dx.
\]
积分在 \([1,2]\) 上没有困难，因为 \(E\) 在该紧区间有界。若 \(K\) 是这个半平面内的紧集，令 \(\delta=\min_{s\in K}(\operatorname{Re}s-\alpha)>0\)。则尾部被常数倍的
\[
 x^{-1-\delta}(\log x)^B
\]
一致控制。对参数求任意固定次数 \(j\) 的导数，只会再出现至多 \((\log x)^j\) 的因子；代换 \(u=\log x\) 后，控制函数的积分化为常数倍的
\[
 \int_{\log2}^\infty e^{-\delta u}u^{B+j}\,\mathrm du<\infty.
\]
故参数积分在紧集上局部一致收敛，且可逐项求导，\(H\) 在 \(\operatorname{Re}s>\alpha\) 全纯。这也说明固定对数幂不会妨碍严格半平面内的延拓。

在原来的 \(\operatorname{Re}s>1\) 中，由\secref{b3:sec:2702}的部分求和表示，
\[
 -\frac{\zeta'(s)}{\zeta(s)}
 =s\int_1^\infty\psi(x)x^{-s-1}\,\mathrm dx
 =\frac{s}{s-1}+H(s).
\]
注意主项为 \(s/(s-1)\)，因为 \(s\int_1^\infty x^{-s}\,\mathrm dx=s/(s-1)\)。它与 \(1/(s-1)\) 只差全纯常数一，但在精确的积分恒等式中应保留这个因子。

左右两侧都是连通半平面 \(\operatorname{Re}s>\alpha\) 上的亚纯函数，并在其非空开子集 \(\operatorname{Re}s>1\) 相等。亚纯函数的恒等定理于是使等式在整个半平面成立。右侧除 \(s=1\) 外全纯。

若 \(\rho\ne1\) 是该半平面内 \(\zeta\) 的 \(m\ge1\) 阶零点，局部可写成
\[
 \zeta(s)=(s-\rho)^m u(s),\quad u(\rho)\ne0,
\]
其中 \(u\) 全纯，并在足够小的圆盘内非零。于是
\[
 -\frac{\zeta'(s)}{\zeta(s)}
 =-\frac m{s-\rho}-\frac{u'(s)}{u(s)}.
\]
其极点留数为非零整数 \(-m\)，不可能可去，和右侧在 \(\rho\) 全纯矛盾。\(s=1\) 已知是极点，不是零点，命题得证。
\end{proof}

这里得到的是 \(\operatorname{Re}s>\alpha\)，没有包括边界 \(\operatorname{Re}s=\alpha\)。在边界上，刚才的控制积分失去了因子 \(e^{-\delta u}\)，原误差假设不再保证收敛。特别地，即使误差是 \(O(x^{1/2}(\log x)^B)\)，这个方向也只直接排除实部严格大于二分之一的零点。

\ProseParagraphBreak

反向从无零点区域推出定量误差，还需要控制 \(\zeta'/\zeta\) 的大小、反演的截断高度和围道边上的积分。仅仅说明某个区域没有零点，还没有完成这些估计。本章的 Tauber 证明按“先固定边界紧段，再取变量趋于无穷”的次序获得 \(o(x)\)；它没有暗含有效的无零区域宽度，也没有自动给出通常更强的误差界。
